% !TEX root = ../main.tex
\begin{figure}[H]
    \centering \includegraphics{01yuadrem/img/30history_i.png}
\end{figure}

\section{History} \label{sec::history}
% History is known in detail thanks to the dutiful oths that recorded it under Tol's guidance.
% TODO: Capitalise all the location names, like Arctic Archipelago, Sylvan Canyon, etc.
% TODO: Take a look at the whole cultures and Yuadrem pages and fix years and nation names!

\subsection*{Ancient History}
\subparagraph{682 BS --- First Communion} In the middle of the dead sea, the et E'ukarilth merges with a deceased higher one embryo.
This transforms the tall one into an insane visage of their former self.
The church of E'ukarilth is later founded to attempt communication with the et.

\subparagraph{592 BS --- Birth of Gats} The search for the Lung of Ur begins, an artifact of great value to the tall kin.
The indigo school of the et Thul-yharch creates the hardy gats, believing the relic is below the surface.

\subparagraph{547 BS --- Birth of Irds} With underground search proving unsuccessful, the red school of Zyl'rech births the mobile irds.
Taking to the skies, they survey land and ocean, hoping to find clues of Lung's location.

\subparagraph{523 BS --- Birth of Marsets} The gold school of Tosh-drieln produces the arboreal marsets.
They explore the thick and dark jungles of Yuadrem with ease.

\subparagraph{451 BS --- Birth of Oths} Under mysterious circumstances, oths are created by the et Tol.
Before disappearing, the tall one teaches them writing, and they begin recording history and compiling the findings of the ets and their progeny with great care and detail.

\subparagraph{397 BS --- Ctereth's Workshop} To cope with the uncontrolled population growth of the new kins, the et Ctereth digs a deep cavern in the middle of the dead sea.
Inside it, the tall one builds a workshop and tirelessly crafts qualar to gift the newborns sentience.

\subsection*{Nadir}
\subparagraph{217 BS --- The Rise of the Spire} The tall kin, apparently done with their search, create the spire at the place where E'ukarilth found the higher one.
They build the stone city of Jan'krug atop the mountain.
The progeny kins, now left alone, are forbidden from accessing the dead sea and, incapable of producing qualar, are forced to fight among themselves.

\subparagraph{209 BS --- First Lost Ones} The first plains gats and chu'ash oths are born, separated from their kins by their lack of qualar.
% While ird and marset lost ones also exist, the lack of a qualars doesn't affect these kins as much as their siblings, perhaps due to their wilder nature.

\subparagraph{179 BS --- First Gat City-states} The gats, always fighting adversity, establish the three city-states of Fiele, Avshen, and Alagyaz.
With careful birth control techniques, they manage to maintain a stable population.

\subparagraph{144 BS --- First Siege} A group of three irds known as ``the feathered sunrise'' infiltrates the dead sea and steal tens of thousands of qualar from Ctereth.
The nations of Krudzal, Harual, and Hulnar are later established by their descendants.

\begin{figure}[H]
    \centering \includegraphics{01yuadrem/img/30history_ii.png}
\end{figure}

\subparagraph{92 BS --- Naenks \& Tsaneks Discovery} Trying to find a home, a group of stray marsets known as the Ovovians, stumble upon the naenks and tsaneks of Drejeck.
These two are inexplicable kins born from mold and fungi respectively.

\subparagraph{51 BS --- First Artificial Qualars} The gat Jirar the bonecarver creates a technique to craft rough qualars imitations.
By passing the practice to the gat's disciples, Jirar unshackles the population number of the kins, and boosts Alagyaz's economy to unprecedented levels.
% To date, only gat master bonecarvers have managed to use the technique. One bonecarver's qualar count usually doesn't go above the thousands, but as populations grow so does the need for qualar.

\subsection*{Great Famine}
\subparagraph{0 --- The Schism} The tall kin's folly causes the schism.
The spire, now revealed to be a dormant volcano, catastrophically erupts.
The event destroys Jan'krug and most of the ets.
The spewed ash blocks off sunlight for four decades, starting the age known as the great famine.

The explosion causes a portal known as the Sizzling Gate to be opened in a cave inside of the spire.
This door leads to Nyx, a strange and primal plane that exists outside of Yuadrem.
From the portal spew forth the foreigner kins: the adventurous tortles, the violent grungs, and the ingenious umans, along with the Nyxborn.

\subparagraph{1 AS --- Second Siege} The foreigner's horde, a great army of tortles, grungs, and umans, siege Ctereth's workshop.
They're successful, and the great number of qualar stolen is used to start their own settlements in Yuadrem.

\subparagraph{4 AS --- Zaloths Discovery} The zaloths, a kin made of fire, ash, thunder, and hail, walk down from the ruins of Jan'krug.
They freely roam Yuadrem, following a nomadic lifestyle that keeps most away from civilized society.

\subsection*{Age of Heroes}
\subparagraph{38 AS --- End of the Great Famine} Satisfied with a death toll in the tens of millions, the ash clouds from the spire disperse, finally ending the great famine.

\subparagraph{57 AS --- Quies Discovery} A group of gat voyagers from Avshen rise up to Jan'krug, finding the city ruined beyond repair, covered by solidified lava.
However, what they do find beneath the ruins are the quies, a new kin.
Quies are the last kin created by the ets, and are brought back to Avshen.
They easily integrate into gat society, despite their physical differences.

\subparagraph{71 AS --- Start of the Eternal War} The newly born kingdom of Krudzal in the north begins a war against the stone giants of the northern territories.
The war rages to this day, with little obtained by the thul'kraka irds.

\subparagraph{99 AS --- Jenkash's Separation} A blossoming nation of qulbaba irds is split into forty-five separate tribes by ideological differences.

\begin{figure}[H]
    \centering \includegraphics{01yuadrem/img/30history_iii.png}
\end{figure}

The tribes that will eventually become Jenkash are bound to constant conflict, unable to establish a unified government for more than a hundred years.

\subparagraph{102 AS --- Third Siege} Inspired by their siblings lost three centuries ago, the army of healing is formed.
Mainly composed of gats and oths, they successfully invade Ctereth's dwellings, then personally bringing the stolen qualar to the bughna gats and the chu'ash oths, re-integrating them into civilized society.

\subparagraph{141 AS --- Birth of Isken} Among the dark forests of the Chirping Wilds, the grung empire of Isken is formed.
Initially secretive, they will soon become one of the most fearsome forces in Yuadrem.

\subparagraph{143 AS --- Babaian Genocide} The grungs of Isken easily crush the marset nation of Baba, systematically killing the marsets until very few are left.

\subparagraph{144 AS --- First Isken-Harual War} Ever hungry for power and land, the Iskean empire attacks the Harualish tribes of the Chirping Wilds.
This is the start of a long sequence of slow and bloody wars that will last for more than two centuries.

\subparagraph{174 AS --- Discovery of the Tides} The oths from the temple of Ignelli, led by Hashim, unearth the phenomenon of the tides, learning of its influence on the kins of Yuadrem.
The discovery revolutionizes the way the kins perceive their own feelings and motivations, and leads to them questioning the nature of sentience itself.

\subparagraph{189 AS --- Fourth Siege} To cope with their ever-growing populations, a temporary alliance is formed between the dratl ird houses of the west and the grung empire of the east.
Their union leads to the fourth and final successful siege of Ctereth, enabling a great growth for the Hulnar and Iskean empires.

\subsection*{Age of Nations}
\subparagraph{195 AS --- Founding of the Seven Kingdoms of the Sea} Ever-growing in numbers, the gat city-states coasting the whaler's sea coalesce into nations, each under its own king.
With all the events happening in one year, the formation of the seven kingdoms of the sea initiate an age of prosperity for the horned and retainer kins.

\subparagraph{201 AS --- Invention of Metal Ships} Edren, a thul'kraka ird from Krudzal, designs and invents the first ironclad ship.
The design, named after the ird's son, Durkin, boosts Krudzal's trading capabilities and kick-starts a great colonization campaign.

\subparagraph{212 AS --- Birth of the Dead Sea Clans} Imitating their neighbors to the north, many uman, dratl ird, and plains gat clans are established in the dead sea.

% These clans however are very different from the civilized kingdoms of the north.
% Warlords are elected by strength, and their territories are as shifting as the erratic sandstorms.

\begin{figure}[H]
    \centering \includegraphics{01yuadrem/img/30history_iv.png}
\end{figure}

\subparagraph{229 AS --- Formation of the Jenkashian Empire} Driven by inner conflict, the irds of the qul archipelago exhaust their natural resources.
This forces them to prematurely end their quarrels, and begin invading and pillaging the surrounding territories.

\subparagraph{231 AS --- Ededian Genocide} The Springwater island is almost completely overtaken by Jenkash, decimating the marset population and forcing most into exile.

\subparagraph{247 AS --- Tidal Sway} Hailing from Ignelli, the oth Narr from the Rashiist school of thought performs an uncanny ritual to harness the power of the tides.
This accidentally triggers the tidal sway.

The oth summons the Sorrow into Yuadrem, ending the life of most Rashiists and ravaging the wildlands entirely, blocking access by land to the southern regions of Yuadrem.

\subparagraph{272 AS --- Fifth Siege} The Iskean grungs, banned from buying artificial qualar from Khedrat, attempt a new siege upon Ctereth's workshop.
This time however they fail, stopped by an unsuspected force: the newly formed dead sea clan of Dzarog.
Dzarog is a clan of umans and gats that live in dens around the spire, and protect Ctereth's caverns for yet unknown reasons.

\subparagraph{281 AS --- Creation of Geomancy} The ird nation of Hairuus, protected from Isken by the splitting mountain range, develop the art of geomancy.
As a test of their mastery of it, they elevate an island at the middle of the shield lake, where their capital, the Nest, is built.

\subparagraph{304 AS --- Invention of Gunpowder} Hailing from the young nation of Sulia, the oth Karmin discovers gunpowder.
With this new firepower, many engineers from Sulia design and build varied weapons, like fire spears, hand-cannons, and muskets.
These new weapons give them a proper combat advantage, allowing them to defend themselves from the savage nomadic tribes of the blank plains, and slowly expand their territories to the east.

\subparagraph{331 AS --- Creation of Windherding} The uncommonly peaceful irds from the Dentrala tribe in Jenkash develop the art of windherding.
The other tribes quickly adapt this art for combat, leading to the Drejeck wars against the naenks of Gannag and the Dratl'fal wars against the declining empire of Hulnar.

\subparagraph{340 AS --- Siszgoel's Independence} Siszgoel, a long-standing colony of Krudzal, declares independence.
The nation of Kaldrathal is born, under the rule of the warrior queen Ialul.
The natural deposits of nitrate in the country's island of residence, Krejek, boosts a powerful gunpowder industry, quickly matching that of Sulia.

\begin{figure}[H]
    \centering \includegraphics{01yuadrem/img/30history_v.png}
\end{figure}

\subparagraph{354 AS --- Birth of Ribinhep} Umans, a kin commonly hunted an enslaved, manage to establish permanent settlement in the isle of rust.
Naming themselves Ribinhep, they start conquering the northern fjords using their unique mercury weapons, fighting under the rule of the frostburn king Kuin.

\subparagraph{389 AS --- End of the Isken-Harual Wars} After 248 years, the Isken-Harual wars end, with Isken crushing almost all of the ird tribes.
The grung empire quickly proceeds to attack the Byurev nation, attempting to conquer territories up north.
They are however stopped by the gats, prepared for such an invasion decades ago.

\subparagraph{411 AS --- Vanishing of Hairuus} The lake-based country of Hairuus suddenly vanishes, soon after elevating new land for their growing capital.
Rumors that the lake is haunted begin spreading, and nations avoid claiming the empty territories and abandoned cities for fear of this mysterious curse.

\subparagraph{440 AS --- Gannag Invasion} Seeing that the Jenkashian forces are focused on conquering the mainland, the armies of Gannag suddenly invades the qul archipelago under the command of Kutsa the sharp.
In few weeks they manage to conquer half of Jenkash's homeland, taking prisoner irds as sacrifices to use as birth corpses.

\subparagraph{461 AS --- Kaldrathal's Conquest} Most of the islands of the arctic archipelago are claimed by Kaldrathal, who establishes a new form of government that tries to represent the taken territories.

\subparagraph{498 AS --- Invention of Blast Weapons} Reut, an engineer from Drer, invents a new use of Sulia's gunpowder: Blast weapons.
Used for close-quarters combat, blast weapons aim to both surprise and immolate the enemy.
Among the most famous examples are the flame vent, the firecrackers and the flaming pole-arms.

\subparagraph{533 AS --- Creation of Wordbinding} In collaboration, the many oth houses of Palegna create the art of wordbinding.
The technique quickly gains traction, as it adds a method for trustless trade between peoples and nations.

% \subparagraph{553 AS --- Na'ane's Founding} A large circle of tsaneks led by Tsehant, tired of their class-based society, made a pilgrimage to the fog gorge.
% They establish in it, and form the independent nation of Na'ane.

\subparagraph{577 AS --- Establishment of Tanethism} The king of Khedrat, Olag the Immortal, establishes the recently born Tanethism as the official religion of the nation.
The other Kingdoms of the Sea follow soon after, and Tanethism is quickly adopted by most gats.
% Here is when bonereading becomes accepted in the seven kingdoms.

\subparagraph{589 AS --- Appearance of Fo} Strange, twisted creatures start attacking any village coasting the shield lake, causing havoc.
Fo, the kinless inhabitant of the nest is quickly blamed for the creation of this creatures, but all attempts to reach the being have failed.

\begin{figure}[H]
    \centering \includegraphics{01yuadrem/img/30history_vi.png}
\end{figure}

\subsection*{Golden Age}
\subparagraph{591 AS --- Hulnar's Demise} The strong alliance between the nations of Khedrat and Sulia defeats Hulnar in the Sylvan wars, allowing both nations to occupy a segment of the Ichor mountains and the entirety of the Sylvan canyon.
This act helps mitigate the pirates' presence in the Whaler's Sea, kick-starting an era of peace and trade for the coastal nations.

\subparagraph{599 AS --- Invention of Steel Firearms} The inventive Kaldrathian engineer Seja combines Krudzal's quench-hardened steel with her new refined gunpowder.
The explosive mix leads to the development of fierce steel-based weapons, including long-range cannons, wheel-lock pistols and sophisticated rifles.

\subparagraph{607 AS --- Invention of the Steam Engine} Away from the economic center of Yuadrem, the Na'anian tsanek Nugut invents the steam engine.
Originally used simply to drain the Na'anian coal mines, the tsaneks were quick to notice its potential and found hundreds of applications for the engine over time.

\subparagraph{621 AS --- The Penance} A surreptitious ritual known only as ``The Penance'' is carried by the citizens of Dzarog.
From the top of the spire, they summon a horrible being known as Cabb Goem-Rlamesh into Yuadrem.
The colossal amalgamate of flesh slowly drags itself towards the east, ferociously protected by the Dzarogian armies.

\subparagraph{628 AS --- Krudzal's First Victory} Using modified Kaldrathal cannons, Krudzal finally manages to kill a stone giant, claiming their first victory in the Eternal War.
% The event strikes fear on the giants, and Krudzal manages to claim their first territories in the mainland.

\subparagraph{635 AS --- Dissolution of Hulnar} Heirless, the king Sul'rech of Hulnar suddenly dies at a young age.
The dwindling kingdom is split into smaller houses, weak ghosts of Hulnar's old glory.

\subparagraph{655 AS --- Cabb Goem-Rlamesh's Launching} The harrowing immensity, Cabb Goem-Rlamesh, reaches the Burnt Ocean and settles some kilometers off the coast of the dry savanna.
% It becomes known as the breathing city.
% All expeditions to the island have ended poorly.

\subparagraph{659 AS --- Separation of Khedrat} The newly-conquered westernmost territories of Khedrat quickly become tired of monarchy, and peacefully claim independence.
Abandoning old traditions, the new countries of Viphogher and Dnomit embrace democracy: a new, king-less form of government.

\subparagraph{671 AS --- Jenkash's Reclaiming} Forced by Gannag to halt their conquest, the Jenkashian empire focused entirely on re-taking the qul archipelago.
Savage battles are fought, and to date they've managed to reclaim most of their lost homeland.

\subparagraph{672 AS --- Present Day}

\newpage
